\unnumberedChapter{The Block Controller}

The previous chapters got us to understand the LCS concepts, to build LCS nodes and power units on different hardware platforms. It all resulted in building an essential piece necessary for any layout, the LCS base station. So far so good. The next big part to address is to build hardware and software to control devices such as signal, turnouts, and so on. Also, we would need to get feedback from the layout where trains exactly and what they currently do. After all, a trains needs to safely get from A to B. Before diving into block controllers, sensors and actors, here is a brief chapter on where the journey will go.

A block controller is an LCS node that can manage one more blocks, each with one or more sections. This part will describe the hardware and firmware developed for the block controller. The block controller features several ports for the blocks and associated turnouts, signals and so on.

// ??? outline what is ahead...


So far, we covered a lot of ground. The Layout Control System rests on the concept of a node with ports and a common bus for nodes to communicate with each other. There are specialized nodes, one of them being the base station responsible for managing the locomotive sessions and generating the DCC signals. The base station is as the name suggests a central piece in a layout. The base station node shown in the previous chapter featured two track outputs, MAIN and PROG. The MAIN track is what powers the layout. However, the base station power module unit allowed for about 2.5 Amps and hence only smaller layouts will just do fine with one base station. For larger layouts, help to power several track sections is needed.

A common approach is to divide the layout in several sections or blocks and power them individually with a booster or block controller. Such a booster or block controller would be very similar to a base station in that it provides control about the track, current consumption limits and so on. In contrast to a base station a booster or block controller would not produce its own DCC signal or manage a locomotive session but rather use the signal distributed via the LCS bus. Being an LCS node, the block controller will offer a rich set of attributes to manage and query the state and configuration. Again, wherever the functionality is the same that can be found on the base station, there will be the same attributes and functions.

But there is more to it. For a safe and automated operation, dividing the layout on to track blocks with block sections, and block signaling concepts are necessary. A train movement is always a planned movement and hence the track route needed for the movement is authorized from a central operations place. The chapter on block signaling control provided a high level overview on the block signaling concepts. This chapter will provide a big step into automating a layout by introducing the block controller.

Conceptually, a booster and a block controller have very much in common. They both handle a section or block, they both need power modules to generate the track power, they both need to offer basic sensor and control capabilities to manage a track section or block. Our concept will just unify these two modules types and we will call them block controller from now on. A block controller is a LCS node that manages a track. There are versions with one channel, i.e. one block, and two channels. Support for track occupancy detection, an essential feature for block control, is optional since simple boosters would not need it if they just manage one block with no sections. The text will just use the term block controller and only mention boosters explicitly when needed. Let's get started with the overall requirements.



